\documentclass[a4paper,10pt]{article}
\usepackage[utf8]{inputenc}
\usepackage[french] {babel}
\usepackage[T1]{fontenc}
\usepackage{lmodern}
\usepackage{graphicx}
\usepackage{graphics}
\usepackage{ulem}
\usepackage{amssymb}
\usepackage{url}
\usepackage[a4paper]{geometry}
\geometry{hscale=0.7,vscale=0.7,centering}
\usepackage{vmargin}
\usepackage{amsmath}
\usepackage{amssymb}
\usepackage{amsthm}
\usepackage{moreverb}
\usepackage{listings}
\newtheorem{theorem}{Théorème}[section]
\newtheorem{defi}{Définition}[section] 
\newtheorem{prop}{Propriété}[section] 
\usepackage{color}
\definecolor{gris}{rgb}{0.95,0.95,0.95}
\lstset{numbers=left, tabsize=4, backgroundcolor=\color{gris},
frame=single, breaklines=true,
keywordstyle=\color{black},
stringstyle=\ttfamily,
framexleftmargin=6mm, xleftmargin=6mm}
%opening
\title{LINGI 2261 : Artificial Intelligence \\
Assignement 2}
\author{Rochet Florentin Debroux Léonard} 
\date{Année académique 2011-2012}

\begin{document}

	\begin{titlepage}
		\begin{center}
			{\huge LINGI2261: Artificial Intelligence}\\
			\vspace{0.4cm}
			
			{\Large {Professor : Yves Deville\\ \vspace{0.2cm} Teaching assistants : Cyrille Dejemeppe and Jean-Baptiste Mairy  }}\\
			\vspace{0.6cm}
			
			{\Large \textit{ Assignement2 : Solving Problems with Informed Search}}\\
			\vspace{1.2cm}

			\texttt{}\\
			\vspace{0.2cm}

			\includegraphics[height=10cm]{pageGarde.png}\\
			\vspace{0.1cm}
			{\Large \textbf{Universit\'e Catholique de Louvain}}
			\vspace{0.7cm}

			Groupe 37 \\
			\vspace{0.2cm}
			
			Florentin Rochet \\
			Léonard Debroux\\
			\vspace{0.2cm}
			2012-2013\\
		\end{center}
	\end{titlepage}

	\newpage
	
	\section{Search Algorithms}
		\subsection{Statments to be proven}
			\subsubsection{Question 1: Prove that depth-first is a special case of uniform-cost search}
				If the tree that is build is such as the cost of the nodes increments by one following the path of an actual depth-search, the uniform-cost will follow that same path and thus perform a depth-first search.
				
			\subsubsection{Question 2: Prove that breadth-first is a special case of uniform-cost search}	
				In the simple case where all the nodes have the exact same cost, the uniform-search will run through the tree by level exactly like a breadth-first search.
			\subsubsection{Question 3: Prove that greedy-best-first search is a special case of A*}
				The f(n) function of an A* search is g(n) + h(n), which is the sum of the exact cost from the beginning to a node n and the extimated cost provided by the heuristic. In a greedy-best first, g(n) is equal to zero and thus it is a special case of A* where we do not take into account the cost to get to the node n. 
			\subsubsection{Question 4: Prove that uniform-cost is a special case of A*}
				This is pretty much the same idea than Q3, only here, it is h(n) that is equal to zero which means that we do not use a heuristic. f(n) is equal only to g(n) that represent the cost from the beginning to the node n.
			
			\subsubsection{Question 5: Prove that if a heuristic is consistent, then it is admissible}
				\begin{proof}
					Let's take the triangle inequality of the consistency and develop it:
					$$ h(n) \leq c(n,a,n') + h(n') $$
					Where node n' is the successor of node n through the action a.\\ 
					Adding g(n), the exact cost to node n, left and right:
					$$ g(n) + h(n) \leq g(n) + c(n,a,n') + h(n') $$
					Simplifying the right part by summing the first two terms:
					$$ g(n) + h(n) \leq g(n') + h(n') $$
					This expression indicates that $$ f(n) = g(n) + h(n) $$ is optimistic and underestimates the cost of a deeper node. Let's say that n' is the goal (and so, h(n)==0), and can be reach by any set of action from node n, we get:
					$$ g(n) + h(n) \leq g(n') $$
					This is the definition of admissibility.
				\end{proof}\\
				We have thus prove that a consistent heuristic is admissible.
				
		\subsection{A* versus uniform-cost search}
			\subsubsection{Question 1: Give a consistent heuristic for this problem. Prove that it is admissible}
				A consistent heuristic for this problem is to use h(n) = manhattan distance from n to the goal. We will prove that it is consistent and thanks to the result of question 1.5, we are sure that it will be admissible too.
				Lets consider a node and the goal, there is two possible cases:\\ 
				Either, we go in the direction of the goal following one of the axes and $$ h(n) = c(n,a,n') + h(n') $$
				Either, we don't move towards the goal and thus we get $$ h(n) < c(n,a,n') + h(n) $$
				In one case of the other, the inequality $$ h(n) \leq c(n,a,n') + h(n) $$ is respect and thus, the heuristic is consistent.\\
				As said earlier, the heuristic is consistent and so, it is admissible.
				
			\subsubsection{Question 2: Show on the left maze the states (board positions) that are visited during an execution of A* graph search with a manhattan distance heuristic (ignoring walls). A state is visited when it is selected in the fringe and expanded. In the A* algorithm, we assume that when different states in the fringe have the smallest value, the algorithm chooses the state with the smallest coordinate (i, j) ((0, 0) being the bottom left position, i being the horizontal index and j the vertical one) using a lexicographical order}
		
			\subsubsection{Show on the right maze the board positions visited by a uniform-cost graph search. When several states have the smallest path cost, this uniform-cost search visits them in the same lexicographical order as A*}
			
	\section{Sokoban planning problem}
	
		\subsection{Question 1: As illustrated on Figure 3 some situations cannot lead to a solution. Are there other similar situations? If yes, describe them.}
			
			There is another situation that cannot lead to a solution, that is when a block is against a wall and must go in a perpendicular direction w.r.t. that wall as shown on the representation below: \\
	\begin{tabular}{ccccc}
		\# & \# & \# & \# & \# \\ 
		\# &    &    &    & \# \\ 
		\# & .  &    & \$ & \# \\ 
		\# &    &    &    & \# \\ 
		\# & \# & \# & \# & \# \\ 
	\end{tabular}\\
	In this cas, the box is stuck against the right wall and it will never be possible to push it to the left to reach the goal.
	In the of a goal being against a wall, there is a dead state if two boxes a pushed against that wall as shown on the exemple below: \\
	\begin{tabular}{ccccccccc}
		\# & \# & \# & \# & \# & \# & \# & \# & \# \\ 
		\# &    & .  &    & \$ &    & \$ &    & \# \\ 
		\# &    &    &    &    &    &    &    & \# \\ 
		\# &    &    &    &    &    &    &    & \# \\ 
		\# & \# & \# & \# & \# & \# & \# & \# & \# \\ 
	\end{tabular}\\
	The box on the right will be able to go on the goal but the one on the right is in the same position as explain above, although there is a goal against the same wall as it.
		
		\subsection{Question 2}
		
		\subsection{Question 3}
		
		\subsection{Question 4}
		
		\subsection{Question 5}
		
		\subsection{Question 6}
		
\end{document}